% !TeX program = xelatex
% ↑ Automatische Auswahl für XeLaTeX compiler

% Das ist mein Template für die TX000 Arbeiten. Nicht perfekt, also falls ihr Verbesserungsvorschläge habt, stellt gerne einen Pull-Request. https://github.com/NikomitK/TX000_Template
% Bitte lasst auch einen star auf github da, danke.
% Wenn ihr die cite funktion von LaTeX nutzen wollt, müsst ihr einfach die Quellen im bibtex Format in die sources.bib datei kopieren, google scholar z.B. hat bei den Quellen einen Button mit dem man das so bekommt, auch viele andere websites.
% Bilder kommen in den images Ordner, den müsst ihr beim abrufen eines Bildes nicht angeben, passiert automatisch.

% Trag hier deine Daten ein, die entsprechenden Felder werden automatisch angepasst.
\def\meinTitel{Title}
\def\meinName{Vorname Nachname}
\def\meinKurs{INF22F}
\def\meineMatrikelNr{Matrikelnummer}
\def\firmenName{Unternehmens-Name}
\def\abteilungsName{Abteilungsname}
\def\projektBetreuer{Vorname Nachname}
\def\ausbildungsLeiter{Vorname Nachname}
\def\abgabeDatum{dd.MM.yyyy}



\documentclass[12pt]{report}
\usepackage[heightrounded]{geometry}
\geometry{
	a4paper,
	lmargin=2.5cm, %Seitenrand left
	tmargin=2.5cm, %Seitenrad top
	headsep=35pt %Abstand von Kopfzeile
}
\usepackage[onehalfspacing]{setspace}
\usepackage[compact]{titlesec}
\usepackage{cite} % Zitierungen
\usepackage{struktex} % Struktogramme
\usepackage{array} % kp hat chatgpt benutzt
\usepackage{longtable} % Tabelle über einen seitenumbruch
\usepackage[nohyperlinks, printonlyused]{acronym} % abkürzungsverzeichnis
\usepackage{fontspec}
\usepackage{blindtext} %LoremIpsum
\usepackage{fancyhdr} %Kopf- und Fußzeile
\usepackage[export]{adjustbox} %Bilder alignment
\usepackage{stfloats} %Tabular at bottom

\usepackage{blindtext}

\usepackage[utf8]{inputenc} % this is needed for umlauts
\usepackage[ngerman]{babel} % this is needed for umlauts
\usepackage[T1]{fontenc}    % this is needed for correct output of umlauts in pdf


\usepackage{subcaption} % Für subfigures glaub


\usepackage{tikz} % Zum zeichnen
\usetikzlibrary{calc}
\usetikzlibrary{shapes.geometric, arrows}
\setmainfont{Arial}

%--------------------Flowcharts--------------------
\tikzstyle{startstop} = [rectangle, rounded corners, minimum width=3cm, minimum height=1cm,text centered, draw=black, fill=red!30]
\tikzstyle{io} = [trapezium, trapezium left angle=70, trapezium right angle=110, minimum width=3cm, minimum height=1cm, text centered, draw=black, fill=blue!30]
\tikzstyle{process} = [rectangle, minimum width=3cm, minimum height=1cm, text centered, draw=black, fill=orange!30]
\tikzstyle{decision} = [diamond, minimum width=3cm, minimum height=1cm, text centered, draw=black, fill=green!30]
\tikzstyle{arrow} = [thick,->,>=stealth]
%--------------------------------------------------


%--------------------Codeblöcke--------------------
\usepackage{listings} %Für Codeblöcke
\usepackage{color} %Farben für Codeblöcke?
\definecolor{dkgreen}{rgb}{0,0.6,0}
\definecolor{gray}{rgb}{0.5,0.5,0.5}
\definecolor{mauve}{rgb}{0.58,0,0.82}

\lstset{
	language=Java,
	aboveskip=3mm,
	belowskip=3mm,
	showstringspaces=false,
	columns=flexible,
	basicstyle={\small\ttfamily},
	numbers=none,
	numberstyle=\tiny\color{gray},
	keywordstyle=\color{blue},
	commentstyle=\color{dkgreen},
	stringstyle=\color{mauve},
	breaklines=true,
	breakatwhitespace=true,
	tabsize=3
}
%-------------------------------------------------

\usepackage{graphicx} %Package für Bilder
\graphicspath{ {./images/} } %Ordner für Bilder

\sloppy % damit lange Wörter nicht über die Zeile hinausgeschrieben werden.

%--------------------Chapter Heading--------------------
\makeatletter
\def\@makechapterhead#1{%
	\vspace*{-20\p@}%
	{\parindent \z@ \raggedright \normalfont
		\ifnum \c@secnumdepth >\m@ne
		%\huge\bfseries \@chapapp\space \thechapter
		\Huge\bfseries \thechapter.\space%
		%\par\nobreak
		%\vskip 20\p@
		\fi
		\interlinepenalty\@M
		\Huge \bfseries #1\par\nobreak
		\vskip 20\p@
}}
\makeatother
\makeatletter
\def\@makeschapterhead#1{%
	\vspace*{-20\p@}%
	{\parindent \z@ \raggedright \normalfont
		%\huge\bfseries \@chapapp\space \thechapter
		\Huge\bfseries\space%
		%\par\nobreak
		%\vskip 20\p@
		\interlinepenalty\@M
		\Huge \bfseries #1\par\nobreak
		\vskip 20\p@
}}
\makeatother
%-------------------------------------------------------

%\addto\captionsngerman{\renewcommand{\listfigurename}{}}
 % titel von abbildungsverzeichnis weg?

\setlength\parindent{0pt} %Auto Einrücken deaktivieren


%-------------Setup für Inhaltsverzeichnis--------------
\renewcommand{\contentsname}{Inhaltsverzeichnis} %Umbenennung TOC

\usepackage{tocloft} % Formatierung TOX

\setlength{\cftbeforetoctitleskip}{0pt}

\renewcommand{\cfttoctitlefont}{\huge\bfseries}
\renewcommand{\cftloftitlefont}{\huge\bfseries}

\renewcommand\cftchapfont{\Large\bfseries}
\renewcommand\cftchappagefont{\large}

\renewcommand\cftsecfont{\large\bfseries}
\renewcommand\cftsecpagefont{\large}

\renewcommand\cftsubsecfont{\large}
\renewcommand\cftsubsecpagefont{\large}

\renewcommand\cftsubsubsecfont{\normalsize}
\renewcommand\cftsubsubsecpagefont{\normalsize}


\renewcommand\cftchapafterpnum{\par\addvspace{8pt}}
\renewcommand\cftsecafterpnum{\par\addvspace{8pt}}
\renewcommand\cftsubsecafterpnum{\par\addvspace{6pt}}
\renewcommand\cftsubsubsecafterpnum{\par\addvspace{6pt}}
%-------------------------------------------------------


%------------------Setup für LoF/LoT--------------------
\makeatletter
\renewcommand{\@cftmakeloftitle}{}
\renewcommand{\@cftmakelottitle}{}
\makeatother
\setlength{\cftfigindent}{0em} % change indentation of e.g. "Figure 1" within list of figures
\renewcommand\cftfigfont{\large}
\renewcommand\cftfigpagefont{\large}
\setlength{\cfttabindent}{0em} % change indentation of e.g. "Figure 1" within list of figures
\renewcommand\cfttabfont{\large}
\renewcommand\cfttabpagefont{\large}
\setlength{\cftbeforeloftitleskip}{0pt}
\setlength{\cftbeforelottitleskip}{0pt}
%-------------------------------------------------------


%----------------Setup für Verlinkungen-----------------
\usepackage{hyperref}
\hypersetup{
	colorlinks,
	citecolor=black,
	filecolor=black,
	linkcolor=black,
	urlcolor=black
}
%-------------------------------------------------------


%-------------Kopf-/Fußzeile für Titlepage--------------
\fancypagestyle{titlepage}
{
	\fancyhead[L]{\includegraphics[scale=0.09]{firmenlogo}}
	\fancyhead[R]{\includegraphics[scale=0.25]{dhbw}}
	\renewcommand{\headrulewidth}{0pt}
	\fancyfoot[C]{}
}
%-------------------------------------------------------

\begin{document} 
	\begin{titlepage}
		\thispagestyle{titlepage}
		\newcommand\HRule{\rule{\textwidth}{1pt}} %Titellinien

		
		\begin{center}
			
			\vspace*{2cm}
			
			%Title
			{ \huge \bfseries \MakeUppercase{\meinTitel}}
			%{ \large \bfseries subTitle}\\[0.4cm]
			
			\vspace*{1.5cm}
			
			%Art der Arbeit
			\Large BACHELORARBEIT // PROJEKTARBEIT // STUDIENARBEIT
			
			\vspace*{3cm}
			
			%Hochschule
			{\LARGE Studiengang Informatik}\\
			{\LARGE an der Dualen Hochschule}\\
			{\LARGE Baden-Württemberg Stuttgart}\\

			\vspace*{2.5cm}
			
			\Large von \meinName
			
			\vspace*{1.5cm}
			
			\Large Abgabedatum: \abgabeDatum

			\begin{table*}[bp]
				\begin{tabular}{l l l l}
					Kurs: & \meinKurs & Matrikelnummer: & \meineMatrikelNr  \\
					Unternehmen: & \firmenName & Ausbildungsleiter: & \ausbildungsLeiter\\
					Abteilung: &  \abteilungsName & Projektbetreuer: &  \projektBetreuer\\
				\end{tabular}
			\end{table*}
			
			
		\end{center}
		
	\end{titlepage}


%------------------Kopf- und Fußzeile-------------------
\spacing{1.5}

\fancypagestyle{plain}{
	\fancyfoot[L]{\meinName\\
		 - \meinKurs}
	\fancyfoot[C]{Seite \thepage\ }% von \pageref{LastPage}}
	\fancyfoot[R]{\abgabeDatum}
}

\pagestyle{fancy}
\fancyhead{}


\fancyhead[L]{\includegraphics[scale=0.09]{firmenlogo}}
\fancyhead[C]{\nouppercase\leftmark}
\fancyhead[R]{\includegraphics[scale=0.25]{dhbw}}

\renewcommand{\footrulewidth}{0.4pt} %Linie für Fußzeile


\renewcommand{\sectionmark}[1]{\markboth{#1}{}} 
%-------------------------------------------------------

\pagenumbering{Roman}
\newpage
\chapter*{Abstract}
\addcontentsline{toc}{chapter}{\protect\numberline{}Abstract}
Hier kommt der Abstract.
\newpage

%------------------Inhaltsverzeichnis-------------------

\addcontentsline{toc}{chapter}{\protect\numberline{}Inhaltsverzeichnis}
\tableofcontents
\addtocontents{toc}{}
\thispagestyle{plain}
%-------------------------------------------------------


\newpage
\chapter*{Ehrenwörtliche Erklärung}
\addcontentsline{toc}{chapter}{\protect\numberline{}Ehrenwörtliche Erklärung}

Ich erkläre hiermit ehrenwörtlich, dass ich die vorliegende Arbeit selbstständig und ohne Benutzung anderer als der angegebenen Hilfsmittel angefertigt habe. Aus den benutzten Quellen, direkt oder indirekt, übernommene Gedanken habe ich als solche kenntlich gemacht.
\\
\\
\\
Diese Arbeit wurde bisher in gleicher oder ähnlicher Form oder auszugsweise noch keiner anderen Prüfungsbehörde vorgelegt und auch nicht veröffentlicht

\vspace{2.5cm}

\parbox{6.5cm}{\hrule\medskip Ort, Datum}  \hfill \parbox{6.5cm}{\hrule\medskip Unterschrift} \hspace{3cm}


%********************************
%Abbildungsverzeichnis
%********************************
\newpage
\chapter*{Abbildungsverzeichnis}
\addcontentsline{toc}{chapter}{\protect\numberline{}Abbildungsverzeichnis}

\listoffigures




%********************************
%Tabellenverzeichnis
%********************************
\newpage
\chapter*{Tabellenverzeichnis}
\addcontentsline{toc}{chapter}{\protect\numberline{}Tabellenverzeichnis}

\listoftables


\newpage
\chapter*{Abkürzungsverzeichnis}
\addcontentsline{toc}{chapter}{\protect\numberline{}Abkürzungsverzeichnis}
\markboth{Abkürzungsverzeichnis}{Abkürzungsverzeichnis} %Benötigt, damit das im header steht

%********************************
%Abkürzungsverzeichnis
%********************************
\begin{acronym}[SOAP]
	\acro{BA}{Beispiel Acronym}
\end{acronym}


%Speichern des page counters, um bei Literaturverzeichnis weiter zu zählen.
\newcounter{frontmatterPage}
\addtocounter{frontmatterPage}{\value{page}} 

\newpage
\pagenumbering{arabic}
\chapter{Einleitung}

%********************************
%Einleitung
%********************************
Dies ist ein Template für TX000 Arbeiten, mit ein paar Beispielen, wie diesem \ac{BA} oder folgendem Bild \ref{fig:example image}

\begin{figure}[h]
	\centering
	\includegraphics[width=0.5 \linewidth]{example}
	\caption{Beispielbild mit einem cuten emoji :3}
	\label{fig:example image}
\end{figure}

Außerdem gibt es noch diese Beispieltabelle am Boden der Seite (durch [b]):

\begin{table}[b]
	\centering
	\begin{tabular}[]{l| c r}
		Spalte 1 & Spalte 2 & Spalte 3 \\
		\hline
		Wert 1 & Wert 2 & Wert 3
	\end{tabular}
	\caption{Beispiel Tabelle}
%	\label{tab:example}
\end{table}

\pagebreak
\section{Beispiel Zitat}
Zitierungen werden mit cite\{Erste Zeile der Quelle in .bib datei\} eingefügt und sehen folgendermaßen aus:\cite{knabke2016grundlagen}

\section{Beispiel Diagramme}

\subsection{Flowchart}
\begin{figure}[h!]
	\centering
	\begin{tikzpicture}[node distance=2cm]
		\node (in1) [process] {Öffnen des Templates};
		\node (pro1) [io, below of=in1] {Dein eigener Input};
		\draw [arrow] (in1) -- (pro1);
		\node (pro2) [process, below of=pro1] {Kompilation};
		\draw [arrow] (pro1) -- (pro2);
		\node (out1) [io, below of=pro2] {Output PDF};
		\draw [arrow] (pro2) -- (out1);
	\end{tikzpicture}
	\caption{Beispiel Flowchart}
	\label{fig:RufRFilter}
\end{figure}

\subsection{Struktogramm}
Hier findet sich die Dokumentation zum Package: \href{https://augenbit.de/wiki/index.php?title=Struktogramme_mit_LaTeX}{Hyperlink beispiel}.
Höhe und breite muss man leider selbst in Millimeter angeben, aber geht schon.

\begin{figure}[h!]
	\begin{struktogramm}(153,55)[]
		\assign{
			\begin{declaration}[Parameter:]
				\description{\pVar{Beispiel Parameter}}{Hier die Parameter der Methode}
			\end{declaration}
		}
		\ifthenelse[12]{3}{3}{DuBenutztDiesesTemplate == true}{\sTrue}{\sFalse}
		
			\return{Du bestehst}
		\change % das \change ist quasi das "else if"
			\forallin{$\forall Formkriterien \in Bewertungsbogen$} % for each schleife
				\sub{Schau mer mal was wird}
			\forallinend
			\return{\sFalse}
		\ifend
	\end{struktogramm}
	\caption{Struktogramm für Bestehen der Praxisarbeit}
	\label{fig:StruktoBestehen}
\end{figure}

\pagebreak
\section{Codeblöcke}
Auch Codeblöcke können dargestellt werden:

\begin{figure}[h!]
	\centering
			\begin{lstlisting}
				public class Main {
					public static void main(String[] args) throws IOException {
						System.out.println("Hallo Welt")
					}
				}
			\end{lstlisting}
	\caption{Beispiel Codeblock}
	\label{fig:BspCodeblock}
\end{figure}


%********************************
%Aufgabenstellung
%********************************
\pagebreak
\section{Aufgabenstellung}
\blindtext
\\
\\
\blindtext


%********************************
%Vorüberlegung
%********************************
\pagebreak
\section{Vorüberlegung}
\blindtext


%********************************
%Lösung
%********************************
\newpage
\chapter{Umsetzung}
\blindtext


%********************************
%Iwas eigenes
%********************************
\pagebreak
\section{Beispiel}
\blindtext


%********************************
%Zusammenfassung und Ausblick
%********************************
\newpage
\chapter{Kritische Reflexion und Ausblick}
\blindtext


%********************************
%Literaturverzeichnis
%********************************
\newpage
\pagenumbering{Roman}
\setcounter{page}{\value{frontmatterPage}} %Bei \pagenumbering wird Seitenzähler zurückgesetzt, hier wird der gespeicherte Wert vom frontmatter weitergeführt
\addtocounter{page}{1}
\addcontentsline{toc}{chapter}{\protect\numberline{}Literaturverzeichnis}

%Quellenverzeichnis
\renewcommand{\refname}{Literaturverzeichnis}
\bibliographystyle{IEEEtran}
\bibliography{./sources}


\end{document}
